\documentclass[../main.tex]{subfiles}

\begin{document}

\chapter*{Abstract}
\label{chap:abstract}
Grasping the overall architecture of large software systems is a challenging task. A plethora of software visualization approaches has been proposed to ease comprehension of the relationships among components and classes in such systems. A core problem is deciding how to place elements and how to connect them meaningfully, avoiding excessive overlaps and intersections that make visualizations cluttered and unintelligible.

This thesis presents a novel approach based on Wave Function Collapse (WFC), a procedural generation algorithm inspired by quantum mechanics, which works on a grid of locally constrained cells, to connect (software) entities while discouraging overlaps and preserving the connectivity imposed by global constraints. We implement this approach in \wassis, a prototype developed in Pharo, and use it to explore the obtained results.

We show how WFC can be tuned to reduce the amount of path crossings and explore its scalability challenges towards representing real-world software systems. Through single- and multi-relationship examples, we evaluate the characteristics of \wassis that produce structured, ``organic'' visualizations, discussing their potential and current limitations.

\end{document}